\section{Conclusion}

We opened with a puzzle: an agent passes 95\% of tests—but did it make 10 strategic root-cause fixes or 100 random modifications? Existing benchmarks can't tell. Our work in controlled settings shows that strategic debugging can be measured via fix-impact-ratio, a metric discriminating agents that target root causes (ratio 3.90: one fix resolves 3-4 test failures) from those using sequential trial-and-error (ratio 1.07: one fix per failure). Experiments reveal strategic debugging operates as a two-stage feedback loop — agents cluster errors by shared characteristics at 2$\times$ random rate (coefficient 1.909, p=0.001), then prioritize fixes yielding 2.15$\times$ higher proportion of high-impact modifications (42.5\% vs 19.8\%) — but transfer to unseen test cases fails (slope ratio 0.82 < 1.5, p=0.504). This clarifies that strategic debugging is a feedback loop effective within revealed test execution, not a learning system enabling autonomous prediction without test execution.

\subsection{Summary}

We addressed the gap between outcome-focused metrics (pass@k) and process-based evaluation by introducing fix-impact-ratio, a framework measuring \emph{how efficiently} agents debug through error feedback. Our main contributions are:

\begin{enumerate}
\item \textbf{Fix-impact-ratio framework with three execution-based metrics} — fix-impact-ratio (tests per modification), clustering coefficient (pattern recognition), held-out slope (generalization attempt) — validated with large effect size (Cohen's d=3.07, p=0.0001) demonstrating strong discriminative power on controlled data.

\item \textbf{Experimental validation of two-stage mechanism} — agents cluster errors by type (coefficient 1.909 vs random 0.950, p=0.001), enabling prioritization of high-impact fixes (42.5\% vs baseline 19.8\%, 2.15$\times$ improvement), explaining how strategic debugging achieves ratio 3.90 vs sequential 1.07. Transfer learning failed (slope ratio 0.82, p=0.504), refuting pattern-based generalization.

\item \textbf{Theoretical clarification} — strategic debugging is a \emph{feedback loop} (error $\rightarrow$ cluster $\rightarrow$ prioritize $\rightarrow$ fix) effective when error messages guide clustering, not a \emph{learning system} (observe $\rightarrow$ extract $\rightarrow$ predict) supporting transfer to unseen tests without execution. Framework measures iterative debugging with test access, complementing pass@k correctness metrics.
\end{enumerate}

\subsection{Future Directions}

This work opens several research directions grounded in our experimental findings:

\textbf{Validating Ecological Validity:} Mock implementation (controlled \texttt{clustering\_strength=0.5}) established metric sensitivity — fix-impact-ratio \emph{can} detect strategic behavior when engineered (d=3.07). Whether production GPT-4 agents exhibit clustering at coefficient > 0.3 rates on real Codeforces problems remains unverified. Immediate extension: replicate h-e1/h-m1/h-m2 with OpenAI API + curated Codeforces subset (solve\_count > 1000, manual quality review) to confirm mock results generalize to real agents and competitive programming datasets.

\textbf{Understanding Transfer Failure:} Pattern memory extracted patterns from revealed test failures but achieved 0\% usage rate on held-out tests. Two competing explanations: (1) pattern quality insufficient (low coverage/precision), or (2) transfer fundamentally information-limited (can't predict without test execution). Test with oracle error patterns (human-labeled, high-quality) — if held-out slope ratio > 1.5 with oracle patterns, transfer limitation is pattern extraction quality; if ratio remains < 1.5 even with perfect patterns, predictive debugging without execution is fundamentally constrained.

\textbf{Extending to Real-World Scenarios:} Framework applies to multi-test scenarios (competitive programming, GitHub issue solving when test suites exist). SWE-bench evaluates issue resolution rates; fix-impact-ratio could measure debugging efficiency \emph{within} resolution process, revealing \emph{how} agents pass repository tests. Extension: analyze SWE-bench patches with test suites, compute fix-impact-ratio and clustering coefficient, compare agent strategies on real production codebases vs synthetic competitive programming.

\textbf{Alternative Mechanism Variants:} Current work tested error-type clustering (syntax/runtime/logic/edge\_case); alternative clustering dimensions worth exploring: semantic similarity of error messages (embeddings-based), code location (group errors by function/module), computational failure mode (timeout, memory error, wrong answer). Each may reveal additional mechanism variants beyond error-type clustering, testing whether two-stage framework generalizes across clustering strategies.

\textbf{Alternative Transfer Mechanisms:} Current pattern memory (extract from revealed errors, apply to held-out tests) failed. Alternative designs worth exploring: few-shot learning (agent sees 2-3 examples of error type + fix, applies to similar held-out tests), meta-learning across problems (learn debugging strategy from Problem 1-40, apply to Problem 41-50), causal code-behavior modeling (simulate test execution without running tests). Each requires different agent capability than post-hoc clustering.

As code generation agents evolve from single-attempt generators to iterative problem-solvers, understanding \emph{how} they improve becomes as critical as \emph{whether} they succeed. Fix-impact-ratio establishes a measurement approach ready for real-world validation, proposing benchmark design beyond pass@k paradigm to measure debugging efficiency — a step toward comprehensive assessment of agentic capability in autonomous software development.
